\section{Jadwal Penelitian}

Penelitian ini direncanakan berlangsung selama dua bulan
(8 minggu) dengan pembagian tahapan yang terstruktur sesuai
dengan fokus utama penelitian, yaitu:

\begin{enumerate}
    \item Analisis stabilitas prediksi multi-langkah (multi-step forecasting)
\end{enumerate}

Penyusunan jadwal dilakukan berdasarkan dependensi antar
komponen sistem, dimulai dari pembangunan lingkungan deterministik,
implementasi model prediktif, hingga evaluasi stabilitas dan
analisis perencanaan berbasis rollout.

\subsection{Rencana Waktu Pelaksanaan}

Rencana waktu pelaksanaan penelitian selama enam bulan
ditunjukkan pada Tabel berikut.

\begin{table}[H]
\centering
\caption{Rencana Jadwal Penelitian (2 Bulan / 8 Minggu)}
\begin{tabular}{|p{5.8cm}|c|c|c|c|c|c|c|c|}
\toprule
\textbf{Kegiatan} & \textbf{M1} & \textbf{M2} & \textbf{M3} & \textbf{M4} & \textbf{M5} & \textbf{M6} & \textbf{M7} & \textbf{M8} \\
\midrule
Spesifikasi lingkungan IPD & \checkmark & \checkmark & & & & & & \\
\midrule
Desain Populasi Strategi Lawan & \checkmark & \checkmark & & & & & & \\
\midrule
Prosedur simulasi dan generasi data & & \checkmark & \checkmark & & & & & \\
\midrule
Representasi dan persiapan dataset & & & \checkmark & \checkmark & & & & \\
\midrule
Perancangan model prediksi (arsitektur) & & & \checkmark & \checkmark & & & & \\
\midrule
Implementasi protokol pelatihan & & & & \checkmark & \checkmark & & & \\
\midrule
Validasi (Confusion Matrix, NLL, CE) & & & & & & \checkmark & & \\
\midrule
Evaluasi (Confusion Matrix, NLL, CE) & & & & & & \checkmark & & \\
\midrule
Analisis eksperimen dan komparasi model & & & & & & & \checkmark & \checkmark \\
\midrule
Penulisan dan revisi laporan & \checkmark & \checkmark & \checkmark & \checkmark & \checkmark & \checkmark & \checkmark & \checkmark \\
\bottomrule
\end{tabular}
\end{table}

\subsection{Dependensi dan Risiko}

Beberapa dependensi penting dalam penelitian ini:

\begin{itemize}
    \item Implementasi rollout planning bergantung pada stabilitas model prediktif pada evaluasi open-loop.
    \item Evaluasi closed-loop tidak dapat dilakukan sebelum estimator nilai $\hat{Q}$ tervalidasi.
    \item Studi ablasi dilakukan setelah konfigurasi utama sistem stabil.
\end{itemize}

Risiko utama penelitian meliputi:

\begin{itemize}
    \item Akumulasi error prediksi pada horizon panjang
    \item Kompleksitas komputasi akibat $\mathcal{O}(N \cdot n \cdot k \cdot M)$
    \item Ketidakstabilan model pada lawan non-stasioner
\end{itemize}

Strategi mitigasi meliputi pengujian bertahap (offline → open-loop → closed-loop),
monitoring konvergensi estimator, serta pembatasan parameter eksperimen
untuk menjaga efisiensi komputasi.
